\section{Proofs for Section~\ref{sec:trimming} and Proposition~\ref{prop:bangbang}}
\label{app:trimming}

This section proves Proposition~\ref{prop:bangbang} and the results of Section~\ref{sec:trimming}, and with them the
uniform bound \eqref{eq:hausdorff} (Corollary~\ref{cor:uniform}). The witnesses, cuts, kept sets, trimming densities,
the constants $c_f$, $C_f$, $c'_f$, $r_Q$ and the modulus $\Omega$ are those of Section~\ref{sec:trimming}.

\subsection{Structure of the trimming}

\begin{proof}[Proof of Proposition~\ref{prop:bangbang}]
The minimizing difference of two polytopes is unique, and the gradient in the weights is the witness, so each
optimal weight vector solves the bathtub problem of Lemma~\ref{lem:trimmedmean}. The objective is continuous and convex
on a compact convex set, so $\mathcal T_\alpha$ is nonempty, compact and convex.

(i) With $u=\sum_iw_i\phi(X_i)$ and $z=\sum_jv_j\phi(Y_j)$ the objective is $\|u-z\|_\Hh^2$, and $u-z$ ranges over the
convex compact set $\mathcal C_\alpha(P_n)-\mathcal C_\alpha(Q_m)$ (Proposition~\ref{prop:georep}), whose element of
minimal norm is unique and has norm $\widehat D_\alpha$.

(ii) The objective $F$ is convex and differentiable, so $(w^\star,v^\star)\in\mathcal T_\alpha$ if and only if $w^\star$ minimizes $w\mapsto\langle\nabla_wF(w^\star,v^\star),w\rangle$ over the capped simplex, and likewise for $v^\star$. By the reproducing property and part (i), at every minimizer
\[
(\nabla_wF)_i=2(K_{XX}w^\star-K_{XY}v^\star)_i=2\langle\phi(X_i),u^\star-z^\star\rangle_\Hh=2\,\widehat d(X_i),
\]
the last equality because $u^\star-z^\star=\widehat d$ is the same at every minimizer, and Lemma~\ref{lem:trimmedmean}, in its
minimum form, describes the minimizers of this linear function. The statement for $v^\star$ follows from $(\nabla_vF)_j=-2\,\widehat d(Y_j)$ and the maximum form
of the same lemma.

(iii) $\widehat D_\alpha=0$ forces $\widehat d=0$ by (i), that is, $u^\star=z^\star$. Conversely, every feasible pair
with $u=z$ attains the value $0$, which is the minimum.
\end{proof}

The population trimming has the same threshold form \eqref{eq:popdensity}.

\begin{lemma}[Threshold form of the population trimming]
\label{lem:popthreshold}
Let Assumption~\ref{ass:regularity} hold, let $\alpha\in[0,1)$, and suppose that
$\dist(\mathcal C_\alpha(P),\mathcal C_\alpha(Q))>0$. Let $d^\star$ be the element of minimal norm of
$\mathcal C_\alpha(P)-\mathcal C_\alpha(Q)$, write $d^\star=a^\star-b^\star$ with $a^\star=\E_P[g^\star\phi]$ for some
$g^\star\in\mathcal G_\alpha(P)$ and $b^\star\in\mathcal C_\alpha(Q)$, and let $t$ be any $(1-\alpha)$-quantile of
$d^\star(X)$ for $X\sim P$. Then $g^\star=\beta$ $P$-almost surely on $\{d^\star<t\}$ and $g^\star=0$ $P$-almost surely
on $\{d^\star>t\}$.
\end{lemma}

\begin{proof}
We show that $g^\star$ minimizes a linear functional of the density, and then that every minimizer of that
functional has the threshold form. The set $\mathcal C_\alpha(P)-\mathcal C_\alpha(Q)$ is convex and weakly compact, as the image of the weakly compact
product of the two regions (Lemma~\ref{lem:attainment}) under the weakly continuous map $(a,b)\mapsto a-b$, hence
norm-closed; so $d^\star$ exists, and it has a decomposition of the stated form by \eqref{eq:region}. For
$a\in\mathcal C_\alpha(P)$, $a-b^\star$ lies in the same set, so $\langle a-a^\star,d^\star\rangle_\Hh\ge0$, and by the
reproducing property $g^\star$ minimizes $g\mapsto\E_P[gU]$ over $\mathcal G_\alpha(P)$, where $U=d^\star(X)$.

The threshold density is such a minimizer, and every minimizer agrees with it off the cut. Since $\Pp(U<t)\le1-\alpha\le\Pp(U\le t)$, there is
$c\in[0,\beta]$ such that $h=\beta\mathbf 1\{U<t\}+c\mathbf 1\{U=t\}$ lies in $\mathcal G_\alpha(P)$, as in the proof of
Lemma~\ref{lem:cvar}. For every $g\in\mathcal G_\alpha(P)$, $(g-h)(U-t)\ge0$ pointwise, because $h=\beta\ge g$ where $U<t$
and $h=0\le g$ where $U>t$, and $\E[(g-h)(U-t)]=\E[gU]-\E[hU]$ since $\E g=\E h=1$. So $h$ is a minimizer, and $(g^\star-h)(U-t)\ge0$ has
expectation at most $0$, hence vanishes almost surely: $g^\star=\beta$ where $U<t$ and $g^\star=0$ where $U>t$.
\end{proof}

\subsection{The parametric rate}

Throughout this subsection $\alpha\in(0,1)$, $\mathcal A=\mathcal C_\alpha(P)$, $\mathcal B=\mathcal C_\alpha(Q)$,
$\mathcal A_n,\mathcal B_m$ are the empirical regions, $\mathrm{sep}(f)=\mathrm{sep}(f;\mathcal A,\mathcal B)$ and
$\mathrm{sep}_n(f)=\mathrm{sep}(f;\mathcal A_n,\mathcal B_m)$ as in \eqref{eq:sep}, and $\xi(f)=\mathrm{sep}_n(f)-
\mathrm{sep}(f)$ is the fluctuation of the separation. For a probability measure $R$ and a bounded measurable $f$ write
\begin{equation}
\begin{aligned}
L_R(f)&=\inf_{g\in\mathcal G_\alpha(R)}\E_R[gf]=\sup_{t\in\R}\big\{t-\beta\,\E_R(t-f)_+\big\},\\
U_R(f)&=\sup_{g\in\mathcal G_\alpha(R)}\E_R[gf]=\inf_{u\in\R}\big\{u+\beta\,\E_R(f-u)_+\big\},
\end{aligned}
\label{eq:twotail}
\end{equation}
the equalities being Lemma~\ref{lem:cvar} applied to $-f$ and to $f$. For $f\in\Hh$ these are
$L_R(f)=-h_{\mathcal C_\alpha(R)}(-f)$ and $U_R(f)=h_{\mathcal C_\alpha(R)}(f)$, with the support functions of
Section~\ref{sec:regions}. The function maximized in $L_R$ is concave with
right and left derivatives $1-\beta R(f\le t)$ and $1-\beta R(f<t)$, so its maximizers are the $t$ with
$R(f<t)\le1-\alpha\le R(f\le t)$, the $(1-\alpha)$-quantiles of $f$ under $R$; likewise the minimizers in $U_R$ are the
$\alpha$-quantiles. By Proposition~\ref{prop:georep}, $\mathrm{sep}(f)=L_P(f)-U_Q(f)$ and
$\mathrm{sep}_n(f)=L_{P_n}(f)-U_{Q_m}(f)$.

\begin{lemma}[A quadratic margin on the unit sphere]
\label{lem:sepmargin}
Let $K\subset\Hh$ be nonempty, convex and closed, write $\mathrm{sep}_K(f)=\inf_{z\in K}\langle f,z\rangle_\Hh$, let
$d^\star$ be its element of minimal norm, suppose $D=\|d^\star\|_\Hh>0$, and put $f^\star=d^\star/D$. Then
$\mathrm{sep}_K(f)\le\langle f,d^\star\rangle_\Hh$ for every $f\in\Hh$, $\mathrm{sep}_K(f^\star)=D=\max_{f\in B_\Hh}
\mathrm{sep}_K(f)$, and
\[
\mathrm{sep}_K(f^\star)-\mathrm{sep}_K(u)\ \ge\ \tfrac D2\,\|u-f^\star\|_\Hh^2
\qquad\text{for every }u\in\Hh\text{ with }\|u\|_\Hh=1 .
\]
For $d'\in\Hh$ with $D'=\|d'\|_\Hh>0$ and $f'=d'/D'$,
$\|d'-d^\star\|_\Hh^2=(D'-D)^2+D'D\,\|f'-f^\star\|_\Hh^2$.
\end{lemma}

\begin{proof}
Since $d^\star\in K$, $\mathrm{sep}_K(f)\le\langle f,d^\star\rangle_\Hh$. For $z\in K$ and $t\in(0,1]$,
$\|d^\star\|_\Hh^2\le\|d^\star+t(z-d^\star)\|_\Hh^2$ by convexity, and letting $t\downarrow0$ gives $\langle d^\star,z-d^\star\rangle_\Hh\ge0$, that is, $\langle f^\star,z\rangle_\Hh\ge D$ on $K$. So
$\mathrm{sep}_K(f^\star)\ge D=\langle f^\star,d^\star\rangle_\Hh\ge\mathrm{sep}_K(f^\star)$, and for $f\in B_\Hh$,
$\mathrm{sep}_K(f)\le\langle f,d^\star\rangle_\Hh\le D$ by Cauchy--Schwarz. For a unit vector $u$,
$\langle u,f^\star\rangle_\Hh=1-\tfrac12\|u-f^\star\|_\Hh^2$, so $\mathrm{sep}_K(u)\le D\langle u,f^\star\rangle_\Hh=
D-\tfrac D2\|u-f^\star\|_\Hh^2$. Finally $\|D'f'-Df^\star\|_\Hh^2=D'^2+D^2-2D'D\langle f',f^\star\rangle_\Hh$, and
substituting $\langle f',f^\star\rangle_\Hh=1-\tfrac12\|f'-f^\star\|_\Hh^2$ gives the identity.
\end{proof}

\begin{lemma}[A basic inequality at two cuts]
\label{lem:basicineq}
Let (A1) and (A2) of Assumption~\ref{ass:regularity} hold, let $\alpha\in(0,1)$, and suppose $D=D_\alpha(P,Q)>0$ and
$\widehat D=\widehat D_\alpha>0$. Let $f^\star=d^\star/D$ and $\widehat f=\widehat d/\widehat D$, let $t^\star$ be a
$(1-\alpha)$-quantile of $f^\star(X)$ for $X\sim P$ and $u^\star$ an $\alpha$-quantile of $f^\star(Y)$ for $Y\sim Q$,
and let $\widehat t$ be the $(n_{1-\alpha}+1)$-th smallest of the values $\widehat f(X_i)$ and $\widehat u$ the
$(m_{1-\alpha}+1)$-th largest of the values $\widehat f(Y_j)$. Then, for every realization of the samples,
\[
\begin{aligned}
\tfrac D2\|\widehat f-f^\star\|_\Hh^2\ \le\ \xi(\widehat f)-\xi(f^\star)
\ \le\ \beta\Big[&(P-P_n)\big\{(\widehat t-\widehat f)_+-(t^\star-f^\star)_+\big\}\\
&+(Q-Q_m)\big\{(\widehat f-\widehat u)_+-(f^\star-u^\star)_+\big\}\Big].
\end{aligned}
\]
\end{lemma}

\begin{proof}
The population separation loses at $\widehat f$ at most the difference of the fluctuations, and the
quadratic margin turns that loss into a squared distance; the right-hand side then comes from evaluating each trimmed
mean at a cut that is optimal for the other measure. The set $\mathcal A-\mathcal B$ is nonempty, convex and closed (proof of
Lemma~\ref{lem:popthreshold}), with element of minimal norm $d^\star$, and $\mathrm{sep}=\mathrm{sep}_{\mathcal
A-\mathcal B}$; the same holds for the polytope $\mathcal A_n-\mathcal B_m$ and $\widehat d$
(Proposition~\ref{prop:bangbang}(i)). By Lemma~\ref{lem:sepmargin}, $\widehat f$ maximizes $\mathrm{sep}_n$ over
$B_\Hh\ni f^\star$, so $\mathrm{sep}(f^\star)-\mathrm{sep}(\widehat f)=[\mathrm{sep}_n(f^\star)-\mathrm{sep}_n(\widehat
f)]+\xi(\widehat f)-\xi(f^\star)\le\xi(\widehat f)-\xi(f^\star)$, and Lemma~\ref{lem:sepmargin} at $u=\widehat f$ bounds
the left-hand side below by $\tfrac D2\|\widehat f-f^\star\|^2_\Hh$.

At $\widehat f$ we bound the fluctuation from above, using the empirical cuts. Since $P_n(\widehat f<\widehat t)\le n_{1-\alpha}/n\le1-\alpha<(n_{1-\alpha}+1)/n\le
P_n(\widehat f\le\widehat t)$, $\widehat t$ is a $(1-\alpha)$-quantile of $\widehat f$ under $P_n$. By the
characterization after \eqref{eq:twotail}, $L_{P_n}(\widehat f)=\widehat t-\beta P_n(\widehat t-\widehat f)_+$, while
$L_P(\widehat f)\ge\widehat t-\beta P(\widehat t-\widehat f)_+$, so $L_{P_n}(\widehat f)-L_P(\widehat f)\le\beta(P-P_n)
(\widehat t-\widehat f)_+$. Likewise $\widehat u$, the $k$-th smallest of the $\widehat f(Y_j)$ with $k=m-m_{1-\alpha}$, is an $\alpha$-quantile of
$\widehat f$ under $Q_m$, as $\alpha\le k/m<\alpha+1/m$ gives $Q_m(\widehat f<\widehat u)<\alpha\le Q_m(\widehat f\le\widehat u)$. So
$U_{Q_m}(\widehat f)=\widehat u+\beta Q_m(\widehat f-\widehat u)_+$ and $U_Q(\widehat f)\le\widehat u+\beta Q(\widehat
f-\widehat u)_+$. Adding, since $\xi=(L_{P_n}-L_P)+(U_Q-U_{Q_m})$,
$\xi(\widehat f)\le\beta[(P-P_n)(\widehat t-\widehat f)_++(Q-Q_m)(\widehat f-\widehat u)_+]$.

At $f^\star$ we bound it from below in the same way, using the population cuts. Now $t^\star$ maximizes the population objective, so $L_P(f^\star)=t^\star-\beta
P(t^\star-f^\star)_+$, while $L_{P_n}(f^\star)\ge t^\star-\beta P_n(t^\star-f^\star)_+$; and $U_Q(f^\star)=u^\star+\beta
Q(f^\star-u^\star)_+$, while $U_{Q_m}(f^\star)\le u^\star+\beta Q_m(f^\star-u^\star)_+$. Hence
$\xi(f^\star)\ge\beta[(P-P_n)(t^\star-f^\star)_++(Q-Q_m)(f^\star-u^\star)_+]$, and subtracting gives the claim.
\end{proof}

The next lemma uses, as in Section~\ref{app:analytic}, $n_\star=nm/(n+m)$ and $e_k=\E|\varepsilon_1+\dots+
\varepsilon_k|/k\le k^{-1/2}$ for independent Rademacher signs $\varepsilon_i$.

\begin{lemma}[A local modulus for the ramps]
\label{lem:localmod}
Let (A1) and (A2) of Assumption~\ref{ass:regularity} hold, let $f^\star\in\Hh$ and $t^\star,u^\star\in\R$ be fixed, and
let $r,\lambda_1,\lambda_2\ge0$. Let $\mathcal H(r)$ be the set of pairs of functions $h_1=(t-f)_+-(t^\star-f^\star)_+$
and $h_2=(f-u)_+-(f^\star-u^\star)_+$ with $f\in\Hh$, $\|f-f^\star\|_\Hh\le r$, $|t-t^\star|\le\kappa r+\lambda_1$ and
$|u-u^\star|\le\kappa r+\lambda_2$, and put $\Gamma(r)=\sup_{(h_1,h_2)\in\mathcal H(r)}[(P-P_n)h_1+(Q-Q_m)h_2]$. Then $\Gamma(r)$
is measurable, nonnegative and nondecreasing in $r$,
\[
\E \Gamma(r)\le2\big[\kappa r/\sqrt{n_\star}+(\kappa r+\lambda_1)e_n+(\kappa r+\lambda_2)e_m\big],
\]
and for every $\chi>0$, with probability at least $1-e^{-\chi}$,
\[
\Gamma(r)\le\kappa r\,L(\chi)+M(\chi),
\qquad
\begin{aligned}
L(\chi)&=2/\sqrt{n_\star}+2(e_n+e_m)+2\sqrt{2\chi/n_\star},\\
M(\chi)&=2(\lambda_1e_n+\lambda_2e_m)+\sqrt{2\chi\big(\lambda_1^2/n+\lambda_2^2/m\big)} .
\end{aligned}
\]
\end{lemma}

\begin{proof}
Write $f=f^\star+v$, $t=t^\star+a$, $u=u^\star+b$, and for each $x$ let $\varphi^1_x(z)=(t^\star-f^\star(x)+z)_+-(t^\star-
f^\star(x))_+$ and $\varphi^2_x(z)=(f^\star(x)-u^\star+z)_+-(f^\star(x)-u^\star)_+$. Both are $1$-Lipschitz and vanish
at $0$, and $h_1(x)=\varphi^1_x(a-v(x))$, $h_2(y)=\varphi^2_y(v(y)-b)$, so $\Gamma(r)$ is the supremum $\Gamma$ of
Lemma~\ref{lem:hinge} with these maps, radius $r$ and $\rho_i=\kappa r+\lambda_i$. That lemma gives the measurability and
the bound on $\E \Gamma(r)$. The choice $(v,a,b)=0$ gives $\Gamma(r)\ge0$, and the classes increase with $r$. Since
$|\varphi^i_x(z)|\le|z|$, every pair has $|h_1|\le2\kappa r+\lambda_1$ and $|h_2|\le2\kappa r+\lambda_2$, so the tail
bound of the lemma holds with $\Delta_i=2(2\kappa r+\lambda_i)$. With probability at least $1-e^{-\chi}$, by
Minkowski's inequality in $\R^2$,
\[
\begin{aligned}
\Gamma(r)&\le\E \Gamma(r)+\sqrt{2\chi\big[(2\kappa r+\lambda_1)^2/n+(2\kappa r+\lambda_2)^2/m\big]}\\
&\le\E \Gamma(r)+\sqrt{2\chi}\,\big[2\kappa r/\sqrt{n_\star}+(\lambda_1^2/n+\lambda_2^2/m)^{1/2}\big],
\end{aligned}
\]
and adding the bound on $\E \Gamma(r)$ and collecting the terms in $\kappa r$ gives the claim.
\end{proof}

\begin{proof}[Proof of Theorem~\ref{cor:trimrate}]
The basic inequality of Lemma~\ref{lem:basicineq} bounds $\tfrac D2\widehat r^{\,2}$ by a supremum whose mean grows
linearly in $\widehat r$, which forces $\widehat r$ to be of order $(n\wedge m)^{-1/2}$; dyadic shells and Markov's
inequality turn this into a rate, as in the proof of \citet[Theorem~3.2.5]{vandervaart1996weak}. We work with the
direction $f^\star=d^\star/D$, so cuts and radii on the scale of $d^\star$ are divided by $D$. Precisely, we show
that for every $\eta\in(0,1)$ there are a constant $C_\eta$, depending on $\eta$ alone, and a size $n_\eta$, depending on
$\eta$, $(P,Q,k,\alpha)$ and $r_Q$, such that for $n\wedge m\ge n_\eta$ each of the four bounds holds with $C_\eta$ in
place of $O_\Pp$ with probability at least $1-\eta$.

We write $\widehat D=\widehat D_\alpha$, $f^\star=d^\star/D$, $\widehat f=\widehat d/\widehat D$ when $\widehat D>0$,
$\widehat r=\|\widehat f-f^\star\|_\Hh$, $t^\star=t_\alpha/D$ and $u^\star=u_\alpha/D$, which are a
$(1-\alpha)$-quantile of $f^\star(X)$ and an $\alpha$-quantile of $f^\star(Y)$, and let $\widehat t$ and $\widehat u$ be
as in Lemma~\ref{lem:basicineq}, so that $\widehat t=\widehat t_n/\widehat D$. Write $\underline c=\min(c_f,c'_f)$ and
$\underline n=n\wedge m$, fix $\eta\in(0,1)$ and put $\eta'=\eta/2$,
\[
\lambda_X=\frac1{c_f}\Big(\varepsilon_n+\frac1n\Big),\qquad
\lambda_Y=\frac1{c'_f}\Big(\varepsilon_m+\frac1m\Big),\qquad
\varepsilon_n=\Big\{\frac{\log(10/\eta')}{2n}\Big\}^{1/2},
\]
with $\varepsilon_m$ defined likewise, and $\tau^{\mathrm S}=\tau^{\mathrm S}_{n,m}(\alpha,\eta'/5)$ as in
\eqref{eq:tauA}. Let $\Gamma$ be the process of Lemma~\ref{lem:localmod} for $f^\star,t^\star,u^\star$,
$\lambda_1=\lambda_X/D$ and $\lambda_2=\lambda_Y/D$. The proof uses four events, each of probability at least
$1-\eta'/5$: the events $E^+$ and $E^-$ of Theorem~\ref{thm:twosided} at level $\eta'/5$, and, by the
Dvoretzky--Kiefer--Wolfowitz inequality with Massart's constant \citep{massart1990dkw},
$E_X=\{\sup_t|F_n(t)-F(t)|\le\varepsilon_n\}$ for the distribution function $F$ of $d^\star(X)$ and the empirical one
$F_n$ of the $d^\star(X_i)$, and the analogous $E_Y$ for the $d^\star(Y_j)$, with $G$, $G_m$ and $\varepsilon_m$. Let
$E=E^-\cap E_X\cap E_Y$, so $\Pp(E^c)\le3\eta/10$. For $n\wedge m$ large enough, $\tau^{\mathrm S}<D$ and $\max(\lambda_X,\lambda_Y)\le r_Q$.

On $E^-$ the problem is nondegenerate: $\widehat D\ge D-\tau^{\mathrm S}>0$. So $\widehat f$ is a unit vector,
$\widehat r\le2$, and Proposition~\ref{prop:bangbang}(ii) applies, with $\widehat A_n=\{\widehat f<\widehat t\}$; also
$A^\star=\{f^\star<t^\star\}$.

The lower density bounds keep the empirical cuts of $f^\star$ near the population ones, within $\lambda_X$ or
$\lambda_Y$, the radius of the Dvoretzky--Kiefer--Wolfowitz event plus one rank, divided by $c_f$ or $c'_f$. By \eqref{eq:qdens}, $F$ is continuous at $t_\alpha$, so $F(t_\alpha)=1-\alpha$, and
integrating the lower density bound gives $F(t_\alpha+\varepsilon)\ge1-\alpha+c_f\varepsilon$ and
$F(t_\alpha-\varepsilon)\le1-\alpha-c_f\varepsilon$ for $0\le\varepsilon\le r_Q$. Let $d^\star_{(n_{1-\alpha}+1)}$ be the
$(n_{1-\alpha}+1)$-th smallest of the $d^\star(X_i)$, so that $F_n(t)\ge(n_{1-\alpha}+1)/n$ exactly when $t\ge
d^\star_{(n_{1-\alpha}+1)}$. On $E_X$, since $c_f\lambda_X=\varepsilon_n+1/n$ and $\lambda_X\le r_Q$,
$F_n(t_\alpha+\lambda_X)\ge1-\alpha+1/n\ge(n_{1-\alpha}+1)/n$ and $F_n(t_\alpha-\lambda_X)\le1-\alpha-1/n<(n_{1-\alpha}+1)/n$,
using $n(1-\alpha)-1<n_{1-\alpha}\le n(1-\alpha)$; so $|d^\star_{(n_{1-\alpha}+1)}-t_\alpha|\le\lambda_X$. On the second sample
let $d^\star_{[k]}$ be the $k$-th smallest of the $d^\star(Y_j)$, $k=m-m_{1-\alpha}$, which is the $(m_{1-\alpha}+1)$-th largest,
with $\alpha\le k/m<\alpha+1/m$. In the same way \eqref{eq:qdensY} gives, on $E_Y$, $G_m(u_\alpha+\lambda_Y)\ge
\alpha+1/m>k/m$ and $G_m(u_\alpha-\lambda_Y)\le\alpha-1/m<k/m$, so $|d^\star_{[k]}-u_\alpha|\le\lambda_Y$. Dividing by
$D$, the corresponding order statistics of the $f^\star(X_i)$ and $f^\star(Y_j)$ are within $\lambda_1$ of $t^\star$ and
within $\lambda_2$ of $u^\star$. Order statistics are $1$-Lipschitz in the sup-norm of the values, and
$|\widehat f(x)-f^\star(x)|\le\kappa\widehat r$ by Cauchy--Schwarz, so
$|\widehat t-t^\star|\le\kappa\widehat r+\lambda_1$ and $|\widehat u-u^\star|\le\kappa\widehat r+\lambda_2$. The pair of
functions in Lemma~\ref{lem:basicineq} therefore belongs to $\mathcal H(\widehat r)$, and the lemma gives
\begin{equation}
\tfrac D2\,\widehat r^{\,2}\le\beta\,\Gamma(\widehat r) .
\label{eq:basicZ}
\end{equation}

A peeling turns \eqref{eq:basicZ} into a rate. We peel the values of $\underline n^{1/2}\widehat r$ into dyadic shells. For
$j\in\mathbb Z$ let $S_j=\{2^{j-1}<\underline n^{1/2}\widehat r\le2^j\}$. On $S_j\cap E$, \eqref{eq:basicZ} and the monotonicity of
$\Gamma$ give $\Gamma(2^j\underline n^{-1/2})\ge D4^j/(8\beta\underline n)$, so $S_j\cap E$ lies in the measurable
event $\{\Gamma(2^j\underline n^{-1/2})\ge D4^j/(8\beta\underline n)\}$. The bound on $\E \Gamma$ in Lemma~\ref{lem:localmod}, with $n_\star^{-1/2}\le(2/\underline n)^{1/2}$ and $e_n,e_m\le\underline n^{-1/2}$, gives
$\underline n\,\E \Gamma(2^j\underline n^{-1/2})\le a2^j+b$ with $a=2(2+\sqrt2)\kappa$, the linear term, and
$b=2\underline n^{1/2}(\lambda_1+\lambda_2)$, the cost
of the cuts, so by Markov's inequality, since
$\Gamma\ge0$, that event has probability at most $8\beta(a2^j+b)/(D4^j)$. Write $\Xi_{\eta'}=\sqrt{\log(10/\eta')/2}+1$.
Since $\underline n\le n$, $\underline n^{1/2}\lambda_X\le\Xi_{\eta'}/c_f$, and likewise for $\lambda_Y$, so
$b\le4\Xi_{\eta'}/(\underline cD)$. Let $j_0$ be the smallest integer with $2^{j_0}\ge\Lambda V_\alpha/D$, for a
number $\Lambda\ge1$ chosen below. The shells $S_j$ with $j>j_0$ cover $\{\underline n^{1/2}\widehat r>2^{j_0}\}$, so
$\{\underline n^{1/2}\widehat r>2^{j_0}\}\cap E$ lies in the union $\mathcal B_{j_0}$ of the events above over $j>j_0$, and summing the geometric
series,
\[
\Pp(\mathcal B_{j_0})\le\frac{8\beta a}{D2^{j_0}}+\frac{8\beta b}{3D4^{j_0}}
\le\frac{8\beta a}{\Lambda V_\alpha}+\frac{32\beta\Xi_{\eta'}}{3\underline c\Lambda^2V_\alpha^2}
\le\frac{16(2+\sqrt2)}\Lambda+\frac{32\Xi_{\eta'}}{3\Lambda^2},
\]
using $V_\alpha\ge\beta\kappa$ and $V_\alpha^2\ge\beta/\underline c$. Choose $\Lambda$, depending only on $\eta$, so
that the right-hand side is at most $\eta/2$. Since $2^{j_0}<2\Lambda V_\alpha/D$, the event
$\{\widehat r>2\Lambda V_\alpha/(D\underline n^{1/2})\}$ lies in $E^c\cup\mathcal B_{j_0}$, of probability at most $3\eta/10+\eta/2<
\eta$ for $n\wedge m$ large, which is the rate of the direction.

The bounds on the witness and the kept set follow from that on the direction. For the witness, Lemma~\ref{lem:sepmargin} gives $\|\widehat d-d^\star\|_\Hh\le|\widehat
D-D|+\sqrt{\widehat DD}\,\widehat r\le\tau^{\mathrm S}+\sqrt2D\widehat r$ on $E^+\cap E^-$ when $\tau^{\mathrm S}<D$, and \eqref{eq:tauA} gives
$\tau^{\mathrm S}_{n,m}(\alpha,\eta'/5)\le\beta\kappa\big(2\sqrt2+4+2\sqrt{\log(5/\eta')}\big)\underline n^{-1/2}$, a multiple of $V_\alpha\underline n^{-1/2}$ depending only on $\eta$. For the kept set, on $E$, if $\widehat f(x)<\widehat t$ and $f^\star(x)\ge t^\star$, then
$f^\star(x)\le\widehat f(x)+\kappa\widehat r<\widehat t+\kappa\widehat r\le t^\star+2\kappa\widehat r+\lambda_1$; if
$f^\star(x)<t^\star$ and $\widehat f(x)\ge\widehat t$, then $f^\star(x)\ge\widehat t-\kappa\widehat r\ge
t^\star-2\kappa\widehat r-\lambda_1$. So $\widehat A_n\triangle A^\star$ lies in the band
$\{|f^\star-t^\star|\le2\kappa\widehat r+\lambda_1\}=\{|d^\star-t_\alpha|\le2\kappa D\widehat r+\lambda_X\}$, and the same computation puts $\{\widehat d=\widehat t_n\}=
\{\widehat f=\widehat t\}$ in it. If its half-width $\ell=2\kappa D\widehat r+\lambda_X$ is at most $r_Q$, the upper density bound of
\eqref{eq:qdens} gives the band mass at most $2C_f\ell$. Finally, $\widehat g_n$ and $g^\star$ take values in
$[0,\beta]$ and agree outside $(\widehat A_n\triangle A^\star)\cup\{\widehat d=\widehat t_n\}\cup\{d^\star=t_\alpha\}$:
both equal $\beta$ on $\widehat A_n\cap A^\star$ and $0$ where $\widehat d>\widehat t_n$ and $d^\star>t_\alpha$. All three
sets lie in the band, which gives the bound on $\int|\widehat g_n-g^\star|\,dP$. On $E$ the half-width is at most $r_Q$ with probability tending to one, and
$\lambda_X\le\Xi_{\eta'}/(c_f\underline n^{1/2})$, while the rate of the direction bounds $2\kappa D\widehat r$ by
$4\Lambda\kappa V_\alpha\underline n^{-1/2}$
outside probability $\eta$.
\end{proof}

\subsection{A sufficient condition for \eqref{eq:qdens}}

The density conditions of Section~\ref{sec:trimming} hold, for instance, under the following condition on the witness.

\begin{remark}[A sufficient condition for \eqref{eq:qdens}]
\label{rem:qdens}
Let $\mathcal X=\R^d$ with the Gaussian kernel, so that $d^\star$ is infinitely differentiable and vanishes at
infinity. If $t_\alpha\neq0$, the level set $\{d^\star=t_\alpha\}$ is nonempty, the gradient of $d^\star$ does not
vanish on it, and $P$ has a Lebesgue density $p$ that is continuous near the level set and positive on it, then
\eqref{eq:qdens} holds, by the coarea formula \citep{evans2015measure}. The same conditions at $u_\alpha$ under $Q$
give \eqref{eq:qdensY} (proof below). For the Laplace kernel on $\R$, by contrast, \eqref{eq:qdens} fails for every
pair and every level $0<\alpha<\alpha^\star$, as the law of $d^\star(X)$ has an atom of mass at least $\alpha$ at the
cut (Proposition~\ref{prop:excess}(iii) and the remark after its proof); whether it can hold for the Laplace kernel in
$\R^d$ is open.
\end{remark}

\begin{proof}[Proof of Remark~\ref{rem:qdens}]
For the Gaussian kernel every function in $\Hh$ is infinitely differentiable \citep[Corollary~4.36]{steinwart2008svm}
and vanishes at infinity \citep[Theorem~4.21 and Lemma~4.23]{steinwart2008svm}, so for $0<\rho<|t_\alpha|/2$ the band
$B_\rho=\{|d^\star-t_\alpha|\le\rho\}$ is compact, and for $\rho$ small enough $\nabla d^\star$ does not vanish and $p$ is
continuous and positive on it, by continuity and compactness. On $B_\rho$ the coarea formula \citep{evans2015measure}
gives $d^\star(X)$ the density $h(t)=\int_{\{d^\star=t\}}p/|\nabla d^\star|\,d\mathcal H^{d-1}$ for $|t-t_\alpha|\le\rho$,
which is continuous, as the flow of $\nabla d^\star/|\nabla d^\star|^2$ carries the level sets smoothly onto one another,
and positive. So \eqref{eq:qdens} holds for $r_Q<\rho$ with $c_f$ and $C_f$ the extremes of $h$ on
$[t_\alpha-r_Q,t_\alpha+r_Q]$, and the same argument for $Q$, $Y$ and $u_\alpha$ gives \eqref{eq:qdensY}.
\end{proof}

\subsection{The trimming at the estimated level}

The regions touch in a single point at $\alpha^\star$, and concentrate in Hausdorff distance $d_H$. Recall the modulus
of contact $\Omega(\epsilon)=\sup\{\|a-\mu_{R^\star}\|_\Hh:a\in\mathcal C_{\alpha^\star}(P),\
\dist(a,\mathcal C_{\alpha^\star}(Q))\le\epsilon\}$.

\begin{lemma}[The contact point]
\label{lem:contact}
Let Assumption~\ref{ass:regularity} hold. Then $\mathcal C_\alpha(P)$ is norm-compact for every $\alpha\in[0,1)$. If
$\alpha^\star<1$, then $R^\star=(P\wedge Q)/(1-\alpha^\star)$ is the only common element of $R_{\alpha^\star}(P)$ and
$R_{\alpha^\star}(Q)$, $\mathcal C_{\alpha^\star}(P)\cap\mathcal C_{\alpha^\star}(Q)=\{\mu_{R^\star}\}$, and
$\Omega(\epsilon)\to0$ as $\epsilon\downarrow0$.
\end{lemma}

\begin{proof}
Norm-compactness is part of Lemma~\ref{lem:attainment}.

The common trimming is unique by the maximality of $P\wedge Q$, and the common embedding by (A3). If $R\in R_{\alpha^\star}(P)\cap R_{\alpha^\star}(Q)$, then $(1-\alpha^\star)R\le P\wedge Q$
(Section~\ref{app:regions}), and both have mass $1-\alpha^\star$, so $R=R^\star$; conversely $R^\star$ lies in both. If $a$ lies in both regions, then
$a=\E_P[g\phi]=\E_Q[h\phi]$ with $g\in\mathcal G_{\alpha^\star}(P)$ and $h\in\mathcal G_{\alpha^\star}(Q)$, the
probability measures $gP$ and $hQ$ have the same embedding, and they are equal by (A3). So $a$ is the embedding of a
common trimming, which is $R^\star$.

That the modulus tends to zero is a compactness argument. If $\Omega(\epsilon)\not\to0$, there are $\epsilon_j\to0$, $c>0$ and
$a_j\in\mathcal C_{\alpha^\star}(P)$ with $\dist(a_j,\mathcal C_{\alpha^\star}(Q))\le\epsilon_j$ and
$\|a_j-\mu_{R^\star}\|_\Hh\ge c$. By compactness a subsequence converges in norm to some
$\bar a\in\mathcal C_{\alpha^\star}(P)$, and $\dist(\bar a,\mathcal C_{\alpha^\star}(Q))=0$, so $\bar a$ lies in the
closed set $\mathcal C_{\alpha^\star}(Q)$ as well. Then $\bar a=\mu_{R^\star}$, which contradicts
$\|\bar a-\mu_{R^\star}\|_\Hh\ge c$.
\end{proof}

Let $\mathcal G=\{x\mapsto(f(x)-t)_+:\ f\in B_\Hh,\ t\in[-\kappa,\kappa]\}$, a class of hinge functions, not to be
confused with the sets $\mathcal G_\alpha(P)$ of trimming densities.

\begin{corollary}[Concentration of the regions, simultaneously in $\alpha$]
\label{cor:uniform}
Under Assumption~\ref{ass:regularity}, for every $n\ge1$ and $\delta\in(0,1)$ there is an event of probability at least
$1-\delta$, namely $\{\sup_{g\in\mathcal G}|(P_n-P)g|\le4\kappa/\sqrt n+\kappa\sqrt{2\log(2/\delta)/n}\}$, which does
not depend on $\alpha$ and on which, for every $\alpha\in[0,1)$,
\begin{equation}
d_H\big(\mathcal C_\alpha(P_n),\mathcal C_\alpha(P)\big)\le r_n(\alpha,\delta)
:=\frac{\kappa}{1-\alpha}\bigg[\frac4{\sqrt n}+\sqrt{\frac{2\log(2/\delta)}n}\,\bigg].
\label{eq:radius}
\end{equation}
\end{corollary}

\begin{proof}
The Hausdorff distance between two convex sets is the largest difference of their support functions, each support
function is an infimum over cuts of means of hinge functions, and so one supremum over the hinge class controls every
level. Write $\mathcal A=\mathcal C_\alpha(P)$ and $\mathcal A_n=\mathcal C_\alpha(P_n)$. For $a\in\mathcal A_n$, \eqref{eq:distmax}
for $\mathcal A$ and $\{a\}$, with $-\langle f,a\rangle_\Hh\le h_{\mathcal A_n}(-f)$, gives
\[
\dist(a,\mathcal A)=\max_{f\in B_\Hh}\{-h_{\mathcal A}(-f)-\langle f,a\rangle_\Hh\}
\le\sup_{f\in B_\Hh}\{h_{\mathcal A_n}(-f)-h_{\mathcal A}(-f)\},
\]
and in the same way
$\dist(a,\mathcal A_n)\le\sup_{f\in B_\Hh}\{h_{\mathcal A}(-f)-h_{\mathcal A_n}(-f)\}$ for $a\in\mathcal A$. As in
the deterministic reduction in the proof of Theorem~\ref{thm:twosided}, Lemmas~\ref{lem:cvar} and~\ref{lem:tcompact} bound
$h_{\mathcal A_n}(-f)-h_{\mathcal A}(-f)$ by $\beta(P_n-P)(-f-t)_+$ at the population minimizer $t\in[-\kappa,\kappa]$, and
$h_{\mathcal A}(-f)-h_{\mathcal A_n}(-f)$ by $\beta(P-P_n)(-f-t_n)_+$ at the empirical one. Since $B_\Hh$ and
$[-\kappa,\kappa]$ are symmetric, these functions lie in $\mathcal G$, so $d_H(\mathcal A_n,\mathcal A)\le\beta
\sup_{g\in\mathcal G}|(P_n-P)g|$, in which only $\beta$ depends on $\alpha$.
Lemma~\ref{lem:hinge} for one sample, with $\varphi=(\cdot)_+$, $r=1$, $\rho_1=\kappa$ and $\Delta_1=2\kappa$ (the functions in
$\mathcal G$ take values in $[0,2\kappa]$), bounds the means of $\sup_g(P_n-P)g$ and $\sup_g(P-P_n)g$ by $4\kappa/\sqrt n$,
and at $e^{-\chi}=\delta/2$ each exceeds its mean by more than $\kappa\sqrt{2\log(2/\delta)/n}$ with probability at most
$\delta/2$.
\end{proof}

\begin{proof}[Proof of Proposition~\ref{prop:alphahat}]
The empirical trimming is nearly in both population regions at $\alpha^\star$, and $\Omega$ says such a
point is close to $\mu_{R^\star}$. That $\Omega(\epsilon)\to0$ is Lemma~\ref{lem:contact}, and one event serves all
levels. Let $E$ be the intersection of the events of
Corollary~\ref{cor:uniform} for the two samples at level $\delta'$, of probability at least $1-2\delta'$, with
$r_n=r_n(\alpha^\star,\delta')$ and $r_m=r_m(\alpha^\star,\delta')$. Let $\hat a=\sum_iw_i\phi(X_i)\in\mathcal C_{\tilde\alpha}(P_n)$ and $\hat b=\sum_jv_j\phi(Y_j)\in\mathcal C_{\tilde\alpha}(Q_m)$,
so that $\|\hat a-\hat b\|_\Hh=\widehat D_{\tilde\alpha}$ by Proposition~\ref{prop:bangbang}(i) and
$\MMD(P_n^w,R^\star)=\|\hat a-\mu_{R^\star}\|_\Hh$.

We pass from the sample regions to the population ones at $\alpha^\star$. On $E$, $\dist(\hat a,\mathcal C_{\tilde\alpha}(P))
\le r_n(\tilde\alpha,\delta')\le r_n$ for every $\tilde\alpha\le\alpha^\star$ simultaneously, and likewise for $\hat b$.
The regions increase with the level and are weakly compact, so there are $a\in\mathcal C_{\alpha^\star}(P)$ with
$\|\hat a-a\|_\Hh\le r_n$ and $b\in\mathcal C_{\alpha^\star}(Q)$ with $\|\hat b-b\|_\Hh\le r_m$. Then $\dist(a,\mathcal C_{\alpha^\star}(Q))\le\|a-b\|_\Hh\le r_n+\widehat D_{\tilde\alpha}+r_m$, by the
triangle inequality through $\hat a$ and $\hat b$, and by the definition of $\Omega$,
$\|a-\mu_{R^\star}\|_\Hh\le\Omega(\widehat D_{\tilde\alpha}+r_n+r_m)$, and the triangle inequality gives the bound.
The second sample is the same argument with $P$ and $Q$ exchanged.
\end{proof}

\begin{proof}[Proof of Corollary~\ref{cor:alphalo}]
Consistency of $\widehat\alpha$ makes the trimmed distance at $\widehat\alpha$ small, as the path falls to zero at most
linearly, and Proposition~\ref{prop:alphahat} turns a small trimmed distance into closeness to $R^\star$. Fix
$\epsilon>0$ and $\delta'\in(0,1)$, and let $E$ be the event of the proof of Proposition~\ref{prop:alphahat} at
level $\delta'/3$, of probability at least $1-2\delta'/3$, with its radii $r_n,r_m$. On $E$, for every level
$\alpha\le\alpha^\star$ the empirical regions are within $r_n$ and $r_m$ of the population ones in Hausdorff distance,
by Corollary~\ref{cor:uniform} and the monotonicity of $r_n(\alpha,\delta')$ in $\alpha$, and the distance between two
sets moves by at most the sum of these, so $\widehat D_\alpha\le D_\alpha+r_n+r_m$. On $E\cap\{\widehat\alpha\le\alpha^\star\}$,
Lemma~\ref{lem:pathconvex}(ii) at $\alpha^\star-\widehat\alpha\in[0,\alpha^\star]$ then gives
\[
\widehat D_{\widehat\alpha}\le\frac{D_0(P,Q)}{\alpha^\star(1-\alpha^\star)}\,
\big(\alpha^\star-\widehat\alpha\big)+r_n+r_m ,
\]
and Proposition~\ref{prop:alphahat} at $\tilde\alpha=\widehat\alpha$, a level at most $\alpha^\star$,
gives $\MMD(P_n^w,R^\star)\le r_n+\Omega(\widehat D_{\widehat\alpha}+r_n+r_m)$ for every
$(w,v)\in\mathcal T_{\widehat\alpha}$, and the same bound for $\MMD(Q_m^v,R^\star)$. Since
$\Omega(t)\to0$ as $t\downarrow0$ and $\Omega$ is nondecreasing, there is $\varepsilon_1>0$ with $\varepsilon_1+\Omega(3\varepsilon_1)\le\epsilon$. By
(A3), $D_0(P,Q)=\MMD(P,Q)>0$ as $\alpha^\star>0$. For $n,m$ large enough, $r_n+r_m\le\varepsilon_1$, and by
Proposition~\ref{prop:loconsistency} the event
$A=\{|\widehat\alpha-\alpha^\star|\le\varepsilon_1\,\alpha^\star(1-\alpha^\star)/D_0(P,Q)\}$ has probability at least
$1-\delta'/3$. On $E\cap A\cap\{\widehat\alpha\le\alpha^\star\}$ both distances are at most
$r_n+\Omega(\varepsilon_1+2(r_n+r_m))\le\varepsilon_1+\Omega(3\varepsilon_1)\le\epsilon$, and by Corollary~\ref{cor:lower} and a union bound this event has
probability at least $1-\delta-\delta'$.
\end{proof}
